\subsection{关于滤子的芽}

设 \(\mathfrak{F}\) 是集合 \(X\) 上的滤子。在集合 \(\mathfrak{P}(X)\)（\(X\) 的所有子集所成之集）上，关系

\[
" \text { there   exists } \quad V \in \mathfrak {F} \text { such   that } \quad M \cap V = N \cap V"
\]

（在 M 与 N 之间）是一个等价关系 R，因为 R 显然是自反且对称的，且若 M、N、P 是 X 的三个子集，满足 \(M \cap V = N \cap V\) 与 \(N \cap W = P \cap W\)，其中 V 与 W 属于 F，则由此可得 \(M \cap (V \cap W) = N \cap (V \cap W) = P \cap (V \cap W)\) 且 \(V \cap W \in F\)，故 R 是传递的。X 的子集 M 模 R 的等价类称为 M 关于 F 的芽；商集 \(\mathfrak{P}(X)/R\) 称为 X 的子集（关于 F）的芽所成之集。

映射 \((\mathbf{M},\mathbf{N})\to \mathbf{M}\cap \mathbf{N}\) 与 \((\mathbf{M},\mathbf{N})\rightarrow \mathbf{M}\cup \mathbf{N}\)（定义在

\[
\mathfrak {P} (\mathbf {X}) \times \mathfrak {P} (\mathbf {X})
\]

上）到 \(\mathfrak{P}(\mathbf{X})\) 中，它们与等价关系 \(\mathbf{R} \times \mathbf{R}\) 和 \(\mathbf{R}\) 相容（《集合论》，R，§ 5，第 8 目）。因为若 \(\mathbf{M} \equiv \mathbf{M}'\) (mod R) 且 \(\mathbf{N} \equiv \mathbf{N}'\) (mod R)，则存在 \(\mathbf{V}\) 与 \(\mathbf{W}\)（\(\mathfrak{F}\) 中的），使得

\[
\begin{array}{r l} & {\mathbf {M} \cap \mathbf {V} = \mathbf {M} ^ {\prime} \cap \mathbf {V} \qquad \mathrm{and} \qquad \mathbf {N} \cap \mathbf {W} = \mathbf {N} ^ {\prime} \cap \mathbf {W},} \\ {\mathrm{so\ that}} & {(\mathbf {M} \cap \mathbf {N}) \cap (\mathbf {V} \cap \mathbf {W}) = (\mathbf {M} ^ {\prime} \cap \mathbf {N} ^ {\prime}) \cap (\mathbf {V} \cap \mathbf {W})} \\ {\mathrm{and}} & {(\mathbf {M} \cup \mathbf {N}) \cap (\mathbf {V} \cap \mathbf {W}) = (\mathbf {M} \cap (\mathbf {V} \cap \mathbf {W})) \cup (\mathbf {N} \cap (\mathbf {V} \cap \mathbf {W}))} \\ & {\qquad = (\mathbf {M} ^ {\prime} \cap (\mathbf {V} \cap \mathbf {W})) \cup (\mathbf {N} ^ {\prime} \cap (\mathbf {V} \cap \mathbf {W}))} \\ & {\qquad = (\mathbf {M} ^ {\prime} \cup \mathbf {N} ^ {\prime}) \cap (\mathbf {V} \cap \mathbf {W}).} \end{array}
\]

过渡到商，这些映射诱导出 \((\mathfrak{P}(\mathbf{X}) / \mathbf{R}) \times (\mathfrak{P}(\mathbf{X}) / \mathbf{R})\) 到 \(\mathfrak{P}(\mathbf{X}) / \mathbf{R}\) 中的两个映射，我们（滥用术语）分别记作 \((\xi, \eta) \to \xi \cap \eta\) 与 \((\xi, \eta) \to \xi \cup \eta\)。直接验证即可确认：关于这些合成法则{[}在整个 \(\mathfrak{P}(\mathbf{X}) / \mathbf{R}\) 上定义的{]}，每个元素都是幂等的，且每个法则都是交换的、结合的，并对另一个法则可分配。此外，关系 \(\xi = \xi \cap \eta\) 与 \(\eta = \xi \cup \eta\) 等价；若（滥用术语）把它们记作 \(\xi \subset \eta\)，则容易验证这个关系是 \(\mathfrak{P}(\mathbf{X}) / \mathbf{R}\) 上的一个序，关于此序 \(\mathfrak{P}(\mathbf{X}) / \mathbf{R}\) 是一个格，它以 \(\varnothing\) 的芽为最小元，以 \(\mathbf{X}\) 的芽为最大元。注意，关系 \(\xi \subset \eta\) 表示存在 \(\mathbf{M} \in \xi\)、\(\mathbf{N} \in \eta\) 与 \(\mathbf{V} \in \mathfrak{F}\)，使得 \(\mathbf{M} \cap \mathbf{V} \subset \mathbf{N} \cap \mathbf{V}\)。

现在设 \(X'\) 是另一个集合，并设 \(\Phi\) 是 F 的某个集合到 \(X'\) 中的所有映射所成之集。\(\Phi\) 上的关系

“存在 \(V \in F\)，使得 f 与 g 在 \(V\) 上都有定义且相等”

（在 \(f\) 与 \(g\) 之间）是一个等价关系 \(\mathbf{S}\)；显然 \(\mathbf{S}\) 是自反且对称的，且若 \(f, g, h\) 是 \(\Phi\) 的三个元素，使得 \(f\) 与 \(g\) 在 \(\mathbf{V} \in \mathfrak{F}\) 上有定义，而 \(g\) 与 \(h\) 在 \(\mathbf{W} \in \mathfrak{F}\) 上有定义且相等，则 \(f\) 与 \(h\) 在 \(\mathbf{V} \cap \mathbf{W} \in \mathfrak{F}\) 上有定义且相等，故 \(\mathbf{S}\) 是传递的。模 \(\mathbf{S}\) 的等价类——这里指映射 \(f\)（集合 \(\mathbf{V} \in \mathfrak{F}\) 到 \(\mathbf{X}'\) 中的）的等价类——称为 \(f\)（关于 \(\mathfrak{F}\)）的芽，而商集 \(\tilde{\Phi} = \Phi / \mathbf{S}\) 称为 \(\mathbf{X}\) 到 \(\mathbf{X}'\) 中的映射（关于 \(\mathfrak{F}\)）的芽所成之集。

注. 1) 每个映射 \(f\)——把子集 \(M\)（\(X\) 的）映入 \(X'\)，其中 \(M\) 属于 \(\mathfrak{F}\)——都模 \(S\) 等价于某个映射 \(f_1\)（\(X\) 到 \(X'\) 中的）（这证明上述术语是合理的）：只需把 \(f\) 延拓到 \(X\) 上，例如在 \(X - M\) 上给它取常值即可。

\begin{enumerate}
\def\labelenumi{\arabic{enumi})}
\setcounter{enumi}{1}
\tightlist
\item
  特征函数 \(\varphi_{\mathbf{M}}\) 与 \(\varphi_{\mathbf{N}}\)——子集 \(\mathbf{M}\) 与 \(\mathbf{N}\)（\(\mathbf{X}\) 的）的特征函数——关于 \(\mathfrak{F}\) 有相同的芽，当且仅当 \(\mathbf{M}\) 与 \(\mathbf{N}\) 关于 \(\mathfrak{F}\) 有相同的芽。
\end{enumerate}

设 \(X''\) 是第三个集合，\(\varphi\) 是 \(X'\) 到 \(X''\) 中的映射，\(\Phi'\) 是 \(\mathfrak{F}\) 的某个集合到 \(X''\) 中的所有映射所成之集。对每个 \(f \in \Phi\)，\(\varphi \circ f\) 属于 \(\Phi'\)；此外，立即可见若 \(g \in \Phi\) 与 \(f\) 关于 \(\mathfrak{F}\) 有相同的芽，则 \(\varphi \circ f\) 与 \(\varphi \circ g\) 关于 \(\mathfrak{F}\) 有相同的芽。

因此这个芽只依赖于芽 \(\tilde{f}\)（\(f\) 关于 \(\mathfrak{F}\) 的），记作 \(\varphi(\tilde{f})\)。这样我们就定义了一个映射（滥用术语，记作 \(\varphi\)），从集合 \(\tilde{\Phi}\)（X 到 \(\mathrm{X}'\) 中的映射的芽所成之集）到集合 \(\tilde{\Phi}'\)（X 到 \(\mathrm{X}''\) 中的映射的芽所成之集）。

现在设 \(\mathbf{X}_i^{\prime}\)（\(1 \leqslant i \leqslant n\)）是一些集合，而

\[
\mathbf {Y} = \prod_ {i = 1} ^ {n} \mathbf {X} _ {i} ^ {\prime}
\]

是它们的积；用 \(\Phi_i\)（相应地，\(\Phi\)）表示 \(\mathfrak{F}\) 的某个集合到 \(\mathbf{X}_i'\)（相应地，\(\mathbf{Y}\)）中的所有映射所成之集。若 \(f_i \in \Phi_i\)（对 \(1 \leqslant i \leqslant n\)），且 \(\mathbf{M}_i \in \mathfrak{F}\) 是 \(f_i\) 的定义域，则映射 \(t \to (f_1(t), \ldots, f_n(t))\) 定义在

\[
\bigcap_ {i = 1} ^ {n} \mathbf {M} _ {i}
\]

上，从而属于 \(\Phi\)；我们（滥用术语）把这个映射记作 \((f_1, \ldots, f_n)\)。此外，若 \(f_i\) 与 \(g_i\) 都属于 \(\Phi_i\) 且关于 \(\mathfrak{F}\) 有相同的芽（对 \(1 \leqslant i \leqslant n\)），则立即可见 \((f_1, \ldots, f_n)\) 与 \((g_1, \ldots, g_n)\) 关于 \(\mathfrak{F}\) 有相同的芽；因此这个芽只依赖于诸芽 \(\tilde{f}_i\)（\(f_i\) 的）。若把它记作 \(\Gamma(\tilde{f}_1, \ldots, \tilde{f}_n)\)，则 \(\Gamma\) 显然是积集

\[
\prod_ {i = 1} ^ {n} \tilde {\Phi} _ {i}
\]

到集合 \(\tilde{\Phi}\) 上的一个双射，其中 \(\tilde{\Phi}_i\)（相应地，\(\tilde{\Phi}\)）表示 X 到 \(\mathbf{X}_i'\)（相应地，Y）中的映射关于 \(\mathfrak{F}\) 的芽所成之集；因此，只要没有混淆的危险，我们将（滥用术语）一般用 \((\tilde{f}_1, \ldots, \tilde{f}_n)\) 代替 \(\Gamma(\tilde{f}_1, \ldots, \tilde{f}_n)\)。

由以上所述，每个映射 \(\psi\)（Y 到集合 \(\mathbf{X}''\) 中的）都定义了映射 \((\tilde{f}_1, \ldots, \tilde{f}_n) \to \psi(\tilde{f}_1, \ldots, \tilde{f}_n)\)

\[
\prod_ {i = 1} ^ {n} \tilde {\Phi} _ {i}
\]

到集合 \(\tilde{\Phi}'\)——X 到 \(\mathbf{X}''\) 中的映射的芽所成之集——中。

特别地，若 \(n = 2\) 且 \(X_1', X_2'\) 与 \(X''\) 都等于同一个集合 \(X'\)（于是 \(\psi\) 是在整个 \(X'\) 上定义的一个合成法则），则 \(\psi\) 诱导出一个在整个集合 \(\tilde{\Phi}\)（\(X\) 到 \(X'\) 中的映射的芽所成之集）上定义的合成法则。容易验证，若 \(X'\) 上给定的法则是结合的（相应地，交换的），则 \(\tilde{\Phi}\) 上相应的法则也是如此；若法则 \(\psi\)（\(X'\) 上的）有单位元 \(e'\)，则关于 \(\mathfrak{F}\) 的常值映射 \(x \to e'\) 的芽是 \(\tilde{\Phi}\) 上相应法则的单位元。最后，若 \(X'\) 上的法则有单位元 \(e'\)，则芽 \(\tilde{f}\)（\(f \in \Phi\) 的）在 \(\tilde{\Phi}\) 中有逆元，当且仅当存在含于 f 的定义域中的 \(V \in F\)，使得 \(f(t)\) 在 \(X'\) 中可逆（对每个 \(t \in V\)）；若对每个 \(t \in V\)，用 \(g(t)\) 表示 \(f(t)\) 的逆元，则 g 的芽 \(\tilde{g}\) 就是 \(\tilde{f}\) 在 \(\tilde{\Phi}\) 中的逆元。特别地，若 \(X'\) 关于法则 \(\psi\) 是一个群，则 \(\tilde{\Phi}\) 关于相应的法则也是一个群；同样，若 \(X'\) 是一个环（相应地，环 A 上的一个代数），则 \(\tilde{\Phi}\) 关于相应的合成法则是一个环（相应地，A 上的一个代数）。

